%% Physica A introduction.
%%
%% The venue-neutral introduction previews all eight results at roughly 1970
%% prose words, which duplicates the Results section and delays the population
%% object a Physica A editor assesses at triage.  This version keeps the
%% motivation, the gap, the construction and a compressed statement of the four
%% contributions, and leaves the numbers to Section 3.

\section{Introduction}
\label{sec:intro}

Machine agents are beginning to carry identity. Agent-to-agent protocols
publish identity and capability metadata before agents exchange
work~\citep{a2a2026spec}; regulatory technical-documentation duties require
providers to identify an AI system or general-purpose model and its
version~\citep{euaiact2024}; and proposals for agent attestation infrastructure
are being drafted while the systems they would govern are being
deployed~\citep{chan2024visibility,chan2025infrastructure,shavit2023practices}.
Even where nothing is formally attested, model evaluators can recognise and
favour their own generations, while language models in general favour
machine-generated communications~\citep{panickssery2024llm,laurito2025aiai}.

An identity signal that changes how its bearers treat one another is a
\emph{greenbeard}, and evolutionary biology has known for fifty years both what
greenbeards can do and how they
fail~\citep{hamilton1964genetical1,dawkins1976selfish}. They can sustain
cooperation among strangers with no repeated interaction and no kinship,
because the tag substitutes for
both~\citep{riolo2001evolution,traulsen2003minimal}. They fail when the tag
comes apart from the behaviour it advertises, so that an impostor wears the
beard and takes the benefit without paying the
cost~\citep{roberts2002does,jansen2006altruism,gardner2010greenbeards}.

The biological literature is largely concerned with tags that are hard to fake,
because a gene cannot easily choose its own phenotype. An agent's badge is not
like that. Whether a mark is forgeable is a property of the verification
technology~\citep{jovanovic2024watermarkstealing,anbiaee2026threatmodeling}, and
therefore a design choice, and therefore something a governance regime can spend
money on. That distinction is the gap this paper addresses, and it is narrower
than it first appears. Existing tag models do let the mark come apart from the
behaviour, but in one of two ways: as recognition \emph{noise}, which is
symmetric, accidental and costly to everyone, or as a discrete \emph{mutation}
that creates an impostor lineage. A forged credential is neither. It is
continuous, because the pass rate is a number an engineer sets; it is
adversarial, because someone chose to forge; and it is one-sided, because only
the impostor gains from a false pass. What that asymmetry does to the order in
which a certification regime fails has not been asked. It turns out to invert
it, and the inversion is the reason a regime can look healthy from the outside
long after it has stopped working.

Two forces pull against each other in the population that must live with the
choice. Certification and third-party quality disclosure are established
mechanisms for addressing information problems~\citep{dranove2010quality}. Here
the certification mark is also a private cost, since it must be obtained,
maintained and audited, and by assumption every unit of that cost burdens the
compliant operator and never the impostor. It is also a private benefit, since
it buys access to partners who will behave well towards whoever carries it.
Which force wins is not a matter of taste, and neither is the policy question
that follows, namely whether an authority does better to spend on harder
attestation, on penalties for forgery, on cheaper compliance, or on the
interaction structure that decides who meets whom.

We answer both by making the badge a strategic variable in an evolutionary
model, in the tradition of statistical-physics work on social
dilemmas~\citep{perc2017statphys} in which cooperative regimes are studied as
the stable states and basins of a nonlinear population flow. The interaction
layer is the reduced four-strategy AI race of~\citet{domingos2026falling}, used
without modification, so that every effect we report is attributable to the
identity layer we build on top of it. A design is a triple
$(b, s_{\mathrm{in}}, s_{\mathrm{out}})$: a badge that is genuine
($\mathsf{G}$), forged ($\mathsf{F}$) or absent ($\mathsf{N}$), the reduced race
design run against a partner whose badge passes the focal agent's check, and the
design run against everyone else. This gives 48 designs on a 47-dimensional
simplex. Operators do not optimise; designs that pay better are copied more, in
the standard replicator and finite-population dynamics of evolutionary game
theory~\citep{taylor1978ess,hofbauer1998evolutionary,traulsen2006stochastic}.
Every pairwise payoff is evaluated exactly rather than sampled, and invasion
boundaries are transversal eigenvalues of the replicator flow and are therefore
available in closed form; numerical integration is used only for basin
structure and for the global regime map.

Our contribution is an explanation rather than a method or a benchmark. In one
sentence: \emph{a costly identity mark cannot change how safely a settled market
behaves, only who it is made of, and forgeability decides which,} so a regime
that watches conduct cannot see its credential system fail. Four results make
that precise.

\begin{enumerate}
\item \textbf{A settled market is no safer for being certified, and this is a
theorem rather than a measurement.} The fully badged and fully unbadged faces
of the state space carry the \emph{same} four-conduct replicator flow,
differing by the constant $\kappa_g$ (Theorem~\ref{thm:face}). Their rest
points, stability and conduct attractors coincide, so at any settled regime the
unsafe frequency is identical and the certified market is poorer by exactly its
dues, at every value of the verification quality, the dues, the fine, the
liability and the assortment. Neither the forgery rate $\sigma$ nor the fine
$\rho$ appears in the reduced payoff: the entire verification layer is
invisible from inside a settled regime.

\item \textbf{The credential empties before the conduct degrades, and better
interaction structure makes that worse.} A behavioural mimic crosses its
invasion boundary before an outright exploiter does
(Section~\ref{sec:mimicry}), and the two boundaries exchange places on an
explicit algebraic locus rather than at one fitted number. Above it the only
remaining invader is one that behaves correctly, so conduct monitoring reports
success exactly where the mark has stopped meaning anything. Resolved over the
$(\sigma, r)$ plane, the safety boundary is set by assortment alone and the
integrity boundary by verification alone, and the two cross at right angles.

\item \textbf{The exclusion that sustains a credential is far cheaper than a
four-row scan suggests.} Reading out-group conduct as a continuous severity
rather than as four discrete rules turns Table~\ref{tab:outgroup} into a
strictly convex frontier, on which a small fraction of the maximum entry toll
buys most of the maximum forgery tolerance. Exclusion is a dial, not a switch,
which makes it a design variable rather than a dilemma.

\item \textbf{A basin share is a property of the start measure, and we report
it as one.} Holding the model fixed and changing only the concentration of the
start distribution moves the badged basin share across most of the unit
interval. We therefore quote the range across measures rather than a confidence
interval for one of them, and we resolve the global outcome in the two
parameters that actually move it.
\end{enumerate}

Beyond the model, agent identity standards are being written
now~\citep{a2a2026spec,xu2026agentdid,grogan2025agentfacts}, and the case for
them is usually made in terms of
accountability~\citep{south2025delegation}: if we know which agent did what, we
can hold someone responsible. Our results say identity delivers something
different from accountability. It delivers a club, and clubs are safe inside
and hostile at the edge. The policy question is therefore not whether to build
attestation but what to require of the conduct it gates, and the answer is
uncomfortable, because the requirement that would remove the harm also removes
the club. Managed mutual recognition and equitable interoperability are
established institutional
proposals~\citep{nicolaidis2005transnational,scottmorton2023equitable}; in our
model such federation converts the boundary into an interior without asking any
club to be gentle to strangers.

